% DO NOT COMPILE THIS FILE DIRECTLY: included by slide.tex.

\begin{frame}[t,plain]
\titlepage
\end{frame}

\begin{frame}{What you will practise}
\begin{itemize}
\item Build a community around a shared operational need.
\item Choose an audience and check the actual sharing path.
\item Encode context that people and tools can use.
\item Turn feedback into corrections and better exports.
\end{itemize}
\note{Audience: CSIRTs, ISAC coordinators and analysts with basic MISP knowledge. Suggested format: 45 minutes of discussion and a 15-minute demonstration. This course focuses on community operation; use the usage and administration courses for interface and deployment details. All examples are fictional.}
\end{frame}

\begin{frame}{Start with a decision, then choose the data}
\begin{itemize}
\item \textbf{Incident response:} find related activity and coordinate help.
\item \textbf{Proactive work:} compare research, sensors and public reports.
\item \textbf{Advanced services:} support forensics and coordinated vulnerability handling.
\item \textbf{Community services:} connect members facing similar threats.
\end{itemize}
\note{Retain the original CSIRT use cases, but treat these as examples, not a formal or exhaustive service taxonomy. Ask participants which decision they cannot make today, who needs to make it and how quickly information must arrive. Raw sensor streams may belong in a log platform; share selected findings with provenance in MISP.}
\end{frame}

\begin{frame}{CIRCL's experience: different communities, different needs}
\begin{itemize}
\item Private-sector and CSIRT communities connect varied capabilities.
\item Sector communities need specialised context and vocabularies.
\item Cross-sector initiatives such as X-ISAC help bridge those groups.
\item Exercise communities provide a place to rehearse collaboration.
\end{itemize}
\note{These are examples drawn from CIRCL's experience in the original course, not a current membership directory. Avoid undated organisation counts or participation percentages. A telecom incident, a bank fraud case and a national CSIRT investigation need different handling, expertise and response times. X-ISAC: https://www.x-isac.org/ .}
\end{frame}

\begin{frame}{Cross-sector sharing needs a purpose}
\begin{itemize}
\item Reused infrastructure or techniques can affect several sectors.
\item Another community may already have analysis or a mitigation.
\item Agree what can cross the boundary before an incident.
\item Share the useful finding with its source and restrictions.
\end{itemize}
\note{A common value or ATT\&CK technique is a lead to investigate, not proof of a common actor. Bridging communities does not imply merging their membership or forwarding all data. Ask who may approve onward sharing and how a recipient can request clarification.}
\end{frame}

\begin{frame}{Write a small community charter}
\begin{itemize}
\item \textbf{Purpose:} priority questions, members and expected benefit.
\item \textbf{People:} coordinators, contributors and escalation contacts.
\item \textbf{Rules:} joining, leaving, handling and permitted use.
\item \textbf{Service:} onboarding, feedback and incident support.
\item \textbf{Review:} agree when to revisit scope and membership.
\end{itemize}
\note{Start with a small trusted group and a concrete use case. Decide who operates the platform, who triages contributions and who resolves a handling mistake. Membership changes require both account and sharing-group review. Source: the MISP community-setup guidelines linked in the reading slide and README.}
\end{frame}

\begin{frame}{Three ways for members to participate}
\centering
\begin{tikzpicture}
\node[bp box,text width=4.0cm] (hub) at (0,0) {Community MISP\\hosted by a coordinator};
\node[bp box] (account) at (-3.75,-3.0) {Hosted member\\user account};
\node[bp box] (own) at (0,-3.0) {Member's own\\MISP instance};
\node[bp box] (sector) at (3.75,-3.0) {Sector hub\\MISP instance};
\draw[bp arrow] (account.north) -- node[bp label,sloped,above] {Web / API access} (hub.south west);
\draw[bp arrow] (own.north) -- node[bp label,right,text width=1.8cm] {Selected exchange} (hub.south);
\draw[bp arrow] (sector.north) -- node[bp label,sloped,above] {Selected exchange} (hub.south east);
\end{tikzpicture}

\par\medskip
\raggedright
Choose according to capacity, ownership and handling needs.
\note{The arrows show possible exchange paths, not automatic bidirectional synchronisation or unrestricted forwarding. A hosted account is access to one instance; the other two models involve separately operated instances. Configure and test each direction, identity, filter and feedback route. A sector hub can reduce operational burden but adds coordination responsibilities.}
\end{frame}

\begin{frame}{Lead by example and lower the first barrier}
\begin{itemize}
\item Publish a small, well-explained example for members to copy.
\item Offer templates, pairing sessions and a named contact.
\item Accept useful partial findings with uncertainty made explicit.
\item Give credit and show how a contribution helped someone.
\end{itemize}
\note{A minimum contribution can be a correction, a question or a sighting. Do not make perfect attribution or a complete incident report a prerequisite for asking for help. Improve quality through review while applying the agreed handling rules from the start.}
\end{frame}

\begin{frame}{Participation is more than publishing events}
\begin{itemize}
\item Analysis, enrichment and requests for assistance all add value.
\item Sightings and false-positive reports improve existing work.
\item Help consumers become confident contributors over time.
\item Avoid volume quotas that reward duplicated or weak data.
\end{itemize}
\note{Replace the old undated 30-percent participation claim and fixed growth timeline with locally measured evidence. Ask what prevents contribution: time, tools, skills, trust or an unresolved handling question. Regular workshops can reveal these barriers. Do not call passive members leeches; consumption can still protect the community.}
\end{frame}

\begin{frame}{Agree handling and privacy before sharing}
\begin{itemize}
\item Identify the purpose, recipients and permission to disclose.
\item Minimise personal, victim and unrelated operational details.
\item Agree retention, correction and incident escalation processes.
\item Resolve legal and contractual questions with the responsible advisers.
\end{itemize}
\note{MISP provides controls, not automatic legal compliance. Applicable GDPR, NIS2, DORA or other obligations depend on the organisation and activity; this course is not a legal determination. Pseudonymisation or delegation does not guarantee anonymity: reports, attachments, times and infrastructure may identify a victim. Compliance references: https://github.com/MISP/misp-compliance .}
\end{frame}

\begin{frame}{Use current TLP 2.0 markings}
\small
\begin{description}
\item[TLP:CLEAR] No TLP limit on onward sharing; other obligations still apply.
\item[TLP:GREEN] Share within the relevant community, away from public channels.
\item[TLP:AMBER] Need-to-know sharing within the recipient organisation and its clients.
\item[TLP:AMBER+STRICT] Need-to-know sharing within the recipient organisation only.
\item[TLP:RED] Individual recipients only; no further disclosure.
\end{description}
\medskip
\textbf{TLP communicates boundaries; configure access separately.}
\note{TLP:CLEAR supersedes TLP:WHITE. The AMBER+STRICT qualifier limits disclosure to the organisation. Use FIRST's full definitions and any additional source restrictions; seek the source's permission before wider disclosure. An ordinary organisation-level MISP audience does not implement a named-person TLP:RED audience. Source: https://www.first.org/tlp/ .}
\end{frame}

\begin{frame}{Align the policy, the audience and the path}
\centering
\begin{tikzpicture}
\node[bp wide] (policy) at (-2.5,0) {Handling policy\\source terms, TLP, purpose};
\node[bp wide] (access) at (2.5,0) {MISP access settings\\distribution, groups, permissions};
\node[bp box,text width=6.0cm] (path) at (0,-2.15) {Review the actual outbound path\\sync, API, feed, file or external service};
\node[bp box,text width=6.0cm] (test) at (0,-4.05) {Test the intended audience and payload\\then exchange and verify};
\draw[bp arrow] (policy.south) -- (path.north west);
\draw[bp arrow] (access.south) -- (path.north east);
\draw[bp arrow] (path) -- (test);
\end{tikzpicture}

\par\medskip
\raggedright
A marking alone does not configure a MISP access rule.
\note{This diagram is a review process, not a representation of a built-in TLP enforcement engine. Tags can drive configured filters or workflows, but those must be reviewed and tested. Access checks combine event and finer-grained distribution, sharing groups and the user's permissions. Administrators and privileged integrations require separate consideration. Source: MISP Book, Synchronization and Sharing.}
\end{frame}

\begin{frame}{Select distribution deliberately}
\begin{itemize}
\item Start sensitive drafts with \textbf{Your organisation only}.
\item Use a \textbf{sharing group} for an explicitly defined set of organisations.
\item Review event, object and attribute distribution together.
\item An attribute cannot widen its event's permitted audience.
\item Test access with both an intended and an excluded recipient.
\end{itemize}
\note{Do not equate a community distribution level with a TLP colour. Attributes can inherit event distribution or be more restrictive; object-level settings also matter. Sharing groups can span instances: inspect organisations, server scope and who may extend membership. Organisation-only is not a defence against privileged local access or an integration acting as that organisation. Consult the sharing documentation for push, pull and internal-instance behaviour rather than assuming a simple hop count.}
\end{frame}

\begin{frame}{Sub-communities need maintained boundaries}
\begin{itemize}
\item Use sharing groups where members can share one platform.
\item Consider separate instances when governance or operations differ.
\item Name an owner for membership and periodic access reviews.
\item Review manual exports, feeds and disconnected transfers too.
\end{itemize}
\note{Separation reduces some mistakes but adds accounts, patching, backups and transfer processes. A feed or air gap does not validate the content or grant disclosure permission. Define what may leave a restricted group and who approves it. Departures require revoking accounts, API keys and memberships; changing access cannot recall copies already received.}
\end{frame}

\begin{frame}{Publication is not a privacy switch}
\begin{itemize}
\item Saving a draft and publishing it are separate actions.
\item Unpublished data can still be visible to authorised users and API clients.
\item Publication can trigger notifications and configured exchange.
\item Preview the destinations, then verify delivery and visibility.
\end{itemize}
\note{Use distribution for confidentiality throughout drafting. Check API consumers and real-time integrations as well as push/pull synchronisation. A successful publish job does not prove every recipient received the update. Unpublishing or tightening distribution cannot guarantee deletion of earlier exports or remote copies. Source: MISP Book, Using the System and Synchronization and Sharing.}
\end{frame}

\begin{frame}{Build a record another analyst can understand}
\begin{itemize}
\item \textbf{Provenance:} source, collection method and original reference.
\item \textbf{Meaning:} what happened, why it matters and what is uncertain.
\item \textbf{Time:} observation window, time zone and review date.
\item \textbf{Action:} suggested use, owner and feedback contact.
\end{itemize}
\note{Keep observations separate from interpretations. first\_seen and last\_seen are observation times, not the event creation or edit timestamp. Leave unknown times unknown. Correlations provide leads; a matching value does not establish common attribution. Use an event report to explain the narrative rather than putting all context into the event title.}
\end{frame}

\begin{frame}{Use the data model to preserve relationships}
\begin{itemize}
\item Group related findings in an event with a meaningful scope.
\item Use typed attributes for individual values.
\item Use object templates for related fields such as an email or file.
\item Link objects and cite the evidence behind the relationship.
\end{itemize}
\note{Avoid one event per isolated value if the values describe one coherent investigation; avoid an ever-growing event containing unrelated investigations. Validate types after importing free text. Keep source UUIDs when exchanging native MISP data where appropriate so updates and corrections can be tracked. The data-model and object-template courses provide worked examples.}
\end{frame}

\begin{frame}{Choose a small, shared vocabulary}
\begin{itemize}
\item Agree taxonomies for handling, confidence and analysis state.
\item Define which tags apply to a whole event or individual attributes.
\item Use galaxy clusters for evidence-backed malware, actors or techniques.
\item Document local conventions; reuse upstream vocabularies first.
\end{itemize}
\note{Confidence, source reliability, severity and handling answer different questions. Avoid contradictory overlapping tags and unsupported actor labels. Local-only tags should not carry handling instructions that a remote recipient needs. If an upstream taxonomy, object template or galaxy needs improvement, contribute the change when permitted; maintain private vocabularies through an approved channel when necessary.}
\end{frame}

\begin{frame}{Enrichment has an outbound data flow}
\begin{itemize}
\item Check what a module sends and which service receives it.
\item Review source terms and handling before submitting a value or sample.
\item Preserve the enrichment source, time and uncertainty.
\item Review results before saving or propagating new claims.
\end{itemize}
\note{A lookup can reveal that an organisation is investigating a domain; a sample can expose victim information. Self-hosted and external modules have different data flows. Do not assume all enrichment is local or that several services repeating one report are independent confirmation. Keep API credentials scoped to the integration and review its permissions.}
\end{frame}

\begin{frame}{One dataset, different decisions}
\centering
\begin{tikzpicture}
\node[bp box,text width=2.55cm] (record) at (0,0) {MISP record\\evidence and context};
\node[bp box,text width=3.6cm] (investigate) at (6.45,1.9) {Investigation\\context, history, relationships};
\node[bp box,text width=3.6cm] (detect) at (6.45,0) {Detection\\selected, relevant indicators};
\node[bp box,text width=3.6cm] (block) at (6.45,-1.9) {Blocking\\stricter review and removal};
\draw[bp arrow] (record.north east) -- node[bp label,sloped,above] {Audience and purpose} (investigate.west);
\draw[bp arrow] (record.east) -- node[bp label,above] {Export policy} (detect.west);
\draw[bp arrow] (record.south east) -- node[bp label,sloped,below] {Impact and freshness} (block.west);
\end{tikzpicture}

\par\medskip
\raggedright
Keep useful context; select action data for each consumer.
\note{All branches must respect permitted use and audience. This is a suggested consumer design, not an automatic MISP pipeline. An investigation can need old or contextual values; detection and blocking have different false-positive costs. Store only information allowed by the retention policy. Source: Best Practices in Threat Intelligence, improving analysis and automated consumers.}
\end{frame}

\begin{frame}{Treat \texttt{to\_ids} as an input to export policy}
\begin{itemize}
\item Enable it for attributes intended for detection use.
\item Leave it off for contextual values such as a report's author.
\item It does not grant access or establish maliciousness.
\item Make export filters explicit; inspect what the consumer receives.
\end{itemize}
\note{Detection exports commonly use the IDS flag; general API queries can return contextual attributes too. A flag alone does not create a blocking decision. Verify the exact export format and query on the deployed version. Typical selection criteria include published state, to\_ids, relevant tags and types, warninglist enforcement and observation age. Source: MISP Book, Using the System and Automation/API.}
\end{frame}

\begin{frame}{Handle false positives without losing the explanation}
\begin{itemize}
\item A warninglist match prompts review of the value's role.
\item Shared hosting or a resolver may be relevant context.
\item Exclude unsuitable values from action exports and explain why.
\item Report the problem to the owner; agree who corrects it.
\end{itemize}
\note{A warninglist hit is not a universal benign verdict, and no hit is not proof of safety. Enable and maintain relevant warninglists; verify that the chosen export enforces them, for example with the supported restSearch enforceWarninglist option. Retaining useful historical context does not mean retaining everything indefinitely. Follow the agreed retention policy. Source: MISP Book, Warning Lists.}
\end{frame}

\begin{frame}{Age changes operational usefulness}
\begin{itemize}
\item Record when the activity was observed, not just imported.
\item Set review intervals appropriate to the data and the action.
\item Use sightings and decay models to support reassessment.
\item Remove stale action data from downstream tools as well.
\end{itemize}
\note{An IP address can change tenant; an old file hash can still help an investigation. Decay scores are model outputs, not truth or automatic deletion. An update timestamp is not a new sighting. Consumers need a reconciliation or removal process; an incremental import that only adds values can leave obsolete blocks in place. See the decaying-indicators course for model configuration.}
\end{frame}

\begin{frame}{Make feedback easy to act on}
\begin{itemize}
\item \textbf{Sightings:} report an observation, false positive or expiration.
\item \textbf{Proposals:} suggest attribute changes for the owner to review.
\item \textbf{Extended events:} add separately owned analysis and context.
\item \textbf{Analyst data:} record notes, opinions and relationships.
\end{itemize}
\note{Sighting types are 0 (sighting), 1 (false positive) and 2 (expiration), not confidence levels. Include observation time and source where appropriate. Review sighting visibility and identity settings: feedback can reveal sensitive activity. Extended events and analyst data have their own distribution considerations. Confirm the configured return path and who triages incoming feedback. Source: MISP Book, Sightings and Analyst Data.}
\end{frame}

\begin{frame}{Close the loop after a correction}
\centering
\begin{tikzpicture}
\node[bp wide] (prepare) at (0,0) {Prepare and review\\source, context, audience};
\node[bp wide] (use) at (5.8,0) {Share and use\\verify the consumer's payload};
\node[bp wide] (feedback) at (5.8,-2.75) {Return feedback\\sightings, proposals, analysis};
\node[bp wide] (correct) at (0,-2.75) {Correct and refresh\\republish, reconcile exports};
\draw[bp arrow] (prepare) -- (use);
\draw[bp arrow] (use) -- (feedback);
\draw[bp arrow] (feedback) -- (correct);
\draw[bp arrow] (correct) -- (prepare);
\end{tikzpicture}

\par\medskip
\raggedright
Assign an owner and a response expectation for each step.
\note{The loop is a process to operate, not a promise that every step is automatic. Correct the source record, republish as appropriate, verify propagation and reconcile downstream exports. A false-positive sighting does not automatically remove a firewall rule. For an urgent mistake, notify affected recipients using the agreed escalation route and track acknowledgement.}
\end{frame}

\begin{frame}{Measure benefit and operational health}
\begin{itemize}
\item How quickly do useful findings and corrections reach members?
\item Which contributions led to a decision or a useful investigation?
\item Are sync jobs, workers and consumer updates completing?
\item Are false positives, stale exports and access issues resolved?
\end{itemize}
\note{Choose a baseline and review period with the community. Counts of events and indicators alone measure volume. Track onboarding success, feedback response and examples of useful collaboration while minimising sensitive member-level reporting. Review patching, backups, scoped API credentials and membership with the platform operator; this is not a replacement for the administration course.}
\end{frame}

\begin{frame}{Demonstration: review a small phishing report}
\begin{itemize}
\item Work in a lab, with organisation-only distribution initially.
\item Create a fictional report containing \texttt{login.example.org}.
\item Add provenance, observation time and a contextual email object.
\item Inspect warninglist hits and decide which fields are detection candidates.
\end{itemize}
\note{Use disposable training data and accounts; do not query or browse the example domain. example.org is reserved for documentation and may match a warninglist: that is deliberate. Do not enable live enrichment, synchronisation or action exports for this exercise. The point is to review structure and intent, not to create a production indicator. Keep contact details and attachments fictional too.}
\end{frame}

\begin{frame}{Exercise: check the audience and the consumer}
\begin{enumerate}
\item Choose a lab sharing group and a compatible handling marking.
\item Check visibility with an included and an excluded lab account.
\item Preview a detection export; explain any missing values.
\item Submit a correction and verify the resulting lab export.
\end{enumerate}
\note{Use a disposable lab without production sync or consumers. Use three organisations: the author, an included recipient and an excluded recipient, with non-administrator test accounts. Expected discussion: the marking alone changes no ACL; an unpublished event is not necessarily invisible; contextual fields should not become blocking inputs; the reserved domain may be excluded by warninglists. Preview the configured published/to\_ids and warninglist behaviour explicitly. Publication, if demonstrated, stays in the isolated lab. Finish by clearing the lab export and checking removal, not just addition.}
\end{frame}

\begin{frame}{Before the next community contribution}
\begin{enumerate}
\item Is the finding useful, sourced and clear about uncertainty?
\item Are the audience, handling and outbound paths appropriate?
\item Are contextual and actionable values distinguished?
\item Can recipients report a problem and receive the correction?
\end{enumerate}
\note{Use this as a lightweight review prompt and improve it with the community. Avoid turning it into a contribution quota or a requirement for perfect intelligence. Invite a participant to explain how their organisation would handle the same finding differently for investigation, detection and blocking.}
\end{frame}

\begin{frame}{References for the operator and trainer}
\small
\begin{itemize}
\item \href{https://github.com/MISP/misp-compliance/blob/master/setting-up-ISACs/guidelines_to_set-up_an_ISAC.md}{Guidelines for setting up an information-sharing community}
\item \href{https://www.misp-project.org/best-practices-in-threat-intelligence.html}{Best Practices in Threat Intelligence}
\item \href{https://www.circl.lu/doc/misp/}{MISP Book: usage, sharing, sightings and warninglists}
\item \href{https://www.first.org/tlp/}{FIRST: TLP 2.0 definitions and usage guidance}
\item \href{https://github.com/MISP/misp-compliance}{MISP compliance resources}
\end{itemize}
\note{Reference review: 3 November 2026. Behaviour depends on the MISP version, permissions, export format and configuration. The course README links individual chapters and records the review decisions. Recheck operational guidance before a live workshop.}
\end{frame}

\begin{frame}{Keep the community conversation going}
\begin{itemize}
\item Share a small useful finding and ask for specific feedback.
\item Bring one handling or integration question to the next meeting.
\item CIRCL: \href{mailto:info@circl.lu}{info@circl.lu} --- \url{https://www.circl.lu/}
\item MISP: \url{https://www.misp-project.org/}
\item Contributions: \url{https://github.com/MISP}
\end{itemize}
\end{frame}
